\slajdid{10-s047}
\begin{frame}[label=10-s047]
\gusttekst\setlength{\abovedisplayskip}{3pt}\setlength{\belowdisplayskip}{3pt}\setlength{\abovedisplayshortskip}{2pt}\setlength{\belowdisplayshortskip}{2pt}\begin{columns}[c,onlytextwidth]\column{.31\textwidth}\centering\begin{tikzpicture}[x=.43cm,y=.58cm,font=\sitneoznake,>=Latex]\draw[->](-.2,0)--(7.1,0)node[right]{$V_I$};\draw[->](0,-.2)--(0,4)node[left]{$V_O$};\draw(0,3.4)--(1.3,3.4)--(3.9,.25)--(6,.25);\node[left]at(0,3.4){$V_{CC}$};\draw[dashed](-.25,.25)node[left]{$V_{CES}$}--(3.9,.25);\draw[dashed](1.3,3.4)--(1.3,-.25);\draw[dashed](3.9,.25)--(3.9,-.65);\node[below]at(1.3,-.25){$(V_D+V_{\gamma T})$};\node[below]at(3.9,-.65){$(2V_D+V_{BES}-V_{\gamma D})$};\node[below]at(6,0){$V_{CC}$};\draw[->,DEplava](6.2,3.1)--(3.3,2.2);\end{tikzpicture}
\column{.67\textwidth}
\[\begin{aligned}V_0&=\frac{V_{CC}-V_{CES}}{(V_D+V_{\gamma T})-(2V_D+V_{BES}-V_{\gamma D})}\\&\quad\cdot\bigl(V_I-(V_D+V_{\gamma T})\bigr)+V_{CC}\end{aligned}\]
\[\begin{aligned}a&=\frac{V_{CC}-V_{CES}}{(V_D+V_{\gamma T})-(2V_D+V_{BES}-V_{\gamma D})}\\&=-\frac{V_{CC}-V_{CES}}{V_D+V_{BES}-V_{\gamma D}-V_{\gamma T}}\end{aligned}\]
\end{columns}\medskip
\begin{columns}[c,onlytextwidth]\column{.39\textwidth}
\[V_I=V_O\]
\[V_I=a\bigl(V_I-(V_D+V_{\gamma T})\bigr)+V_{CC}\]
\[V_I=\frac{V_{CC}-a(V_D+V_{\gamma T})}{(1-a)}=V_M\]
\column{.59\textwidth}
Margine šuma za višestruke izvore šuma su
\[NM_{LMS}=V_{IL}-V_{OL}=V_D+V_{\gamma T}-V_{CES}\]
\[NM_{HMS}=V_{OH}-V_{IH}=V_{CC}-2V_D-V_{BES}+V_{\gamma D}\]
Zbog uske prelazne zone margine šuma za jednostruke izvore će biti slične po vrednosti.
\[NM_{LSS}=V_M-V_{OL}=V_M-V_{CES}\]
\[NM_{HSS}=V_{OH}-V_M=V_{CC}-V_M\]
\end{columns}
\end{frame}
